% ===================================================================
%  DRAFT: The open primitive and the gravitational genus expansion
%
%  This is a draft recast of Vol II foundations, written from the
%  correct primitives. It replaces the discovery narrative ("we found
%  Swiss-cheese inside the bar complex") with the dependency narrative
%  ("the Swiss-cheese operad is primitive; the bar complex is a
%  coalgebraic projection; the derived center is the universal bulk;
%  modularity is trace plus clutching on the open sector; and the
%  genus expansion of the MC element IS gravity").
% ===================================================================

\section{The open factorization category}
\label{sec:open-primitive}

\subsection{Statement of the problem}

A chiral algebra on a curve is a factorization algebra on the
Ran space, with the chiral product encoding collisions of
points.  Volume~I constructed the bar complex~$\barB(\cA)$,
proved five theorems about it, and extracted a tower of modular
invariants.  All of this operates in a single color: the
holomorphic direction.

Three-dimensional holomorphic--topological field theory on
$\C_z \times \R_t$ introduces a second color: the topological
direction~$\R$.  Observables factorize holomorphically in~$z$ and
associatively in~$t$.  The governing structure is a two-colored
operad with closed (holomorphic) and open (topological) colors.

The present chapter identifies the correct primitive datum and
derives everything from it.

\subsection{The Swiss-cheese operad}\label{subsec:sc-operad}

\begin{definition}[The holomorphic--topological Swiss-cheese operad]
\label{def:SC-chtop}
\index{Swiss-cheese operad!holomorphic--topological|textbf}
The \emph{holomorphic--topological Swiss-cheese operad}
$\SCchtop$ is the two-colored topological operad with
colors $\{\mathrm{ch},\, \mathrm{top}\}$ defined by:
\begin{enumerate}[label=\textup{(\alph*)}]
\item \emph{Closed-output spaces.}
  $\SCchtop(\mathrm{ch}^{\times k};\, \mathrm{ch})
  := C_*(\FM_k(\C);\, \cL_{\log})$,
  the chains on the Fulton--MacPherson compactification of~$k$
  points in~$\C$ with coefficients in the logarithmic local system.

\item \emph{Open-output spaces.}
  $\SCchtop(\mathrm{ch}^{\times k},\,
  \mathrm{top}^{\times m};\, \mathrm{top})
  := C_*(\FM_k(\C) \times \Conf_m^{\mathrm{ord}}(\R))$,
  the chains on the product of holomorphic configurations and
  ordered topological configurations.

\item \emph{Directionality.}
  $\SCchtop(\ldots;\, \mathrm{ch}) = \varnothing$ whenever any
  input has color $\mathrm{top}$.  No open input can produce a
  closed output.  This is the operadic expression of the
  bulk$\to$boundary directionality: bulk operators restrict to
  boundaries, but boundary operators do not extend to the bulk.
\end{enumerate}
Composition is induced by the operadic compositions of
$\FM_\bullet(\C)$ (nested collision) and $E_1$ (interval
concatenation).  The factorization product decomposition
$\omega_k = \omega_k^{\mathrm{hol}} \otimes
\omega_k^{\mathrm{top}}$ on $\FM_k(\C) \times \Conf_k(\R)$
governs the weight forms.
\end{definition}

\begin{theorem}[Homotopy-Koszulity; {\cite{Liv15,GK94}};
  \ClaimStatusProvedHere]
\label{thm:homotopy-koszul-SC}
The operad $\SCchtop$ is homotopy-Koszul: its bar construction
$B(\SCchtop)$ is formal \textup{(}quasi-isomorphic to its
cohomology, the Koszul dual cooperad\textup{)}.
\end{theorem}

\begin{proof}
The classical Swiss-cheese operad $\mathsf{SC}$ is Koszul
\cite{Liv15,GK94}.  Kontsevich formality gives a quasi-isomorphism
$\Phi \colon \SCchtop \xrightarrow{\sim} \mathsf{SC}$.
Homotopy transfer through~$\Phi$ gives a quasi-isomorphism
$B(\SCchtop) \xrightarrow{\sim} B(\mathsf{SC})$; the latter is
formal by the Koszul property.  Two-out-of-three completes the
proof.
\end{proof}

\begin{corollary}[Quillen equivalence]
\label{cor:quillen-equivalence}
On the chirally Koszul locus, bar-cobar duality for $\SCchtop$
is a Quillen equivalence: the derived unit and counit are both
weak equivalences.  Consequently, the homotopy categories of
$\SCchtop$-algebras and of coalgebras over the Koszul dual
cooperad are equivalent.
\end{corollary}

\subsection{The open factorization category}
\label{subsec:open-factorization-category}

\begin{definition}[Tangential log curve]
\label{def:tangential-log-curve-v2}
A \emph{tangential log curve} is a triple $(X, D, \tau)$, where
$X$ is a smooth algebraic curve, $D = \{p_1, \ldots, p_n\}$ is a
finite set of marked points, and
$\tau = \{\tau_{p_i} \in T_{p_i} X \setminus \{0\}\}$ is a choice
of nonzero tangent direction at each puncture.
The real oriented blowup $\widetilde{X}_D \to X$ replaces each
$p_i$ by the boundary circle $S^1_{p_i}$; the tangent direction
$\tau_{p_i}$ selects a basepoint, cutting the circle into a
boundary interval $I_{p_i} \cong \R$.
\end{definition}

\begin{definition}[Open/closed factorization category]
\label{def:oc-fact-cat}
\index{open/closed factorization category|textbf}
Let $(X, D, \tau)$ be a tangential log curve.  An \emph{open/closed
factorization dg-category} on $(X, D, \tau)$ is a $\SCchtop$-monoidal
dg-category $\cC$ equipped with:
\begin{enumerate}[label=\textup{(\roman*)}]
\item \emph{Bulk sector.}  A factorization algebra $\cA$ on
  $\Ran(X \setminus D)$ (the closed color).

\item \emph{Boundary sector.}  For each puncture $p_i$, a
  dg-category $\cC(I_{p_i})$ of boundary conditions on the
  interval~$I_{p_i}$ (the open color), with associative
  composition from $E_1$-factorization on~$\R$.

\item \emph{Mixed operations.}  For each configuration of $k$
  bulk points in $X \setminus D$ and $m$ boundary points on the
  intervals~$I_{p_i}$, a mixed operation
  $\mu_{k;m} \colon \cA^{\otimes k} \otimes
  \cC^{\otimes m} \to \cC$
  governed by the mixed-output component of $\SCchtop$.

\item \emph{Compatibility.}  The $\SCchtop$-coherence
  identities (Stokes cancellation on all codimension-one
  boundary strata of $\FM_k(\C) \times \Conf_m(\R)$).

\item \emph{Cyclic trace.}  A nondegenerate trace
  $\Tr_b \colon \RHom_{\cC(I)}(b,b) \to k[-d]$
  for each compact generator~$b$, inducing a Calabi--Yau
  structure on $\cC(I)$.

\item \emph{Clutching.}  For each pair of punctures
  $(p_i, p_j)$, a sewing functor
  $\xi_{ij} \colon \cC(I_{p_i}) \otimes \cC(I_{p_j}) \to
  \cC(I_{p_{ij}})$
  implementing nodal degeneration: the boundary circles $S^1_{p_i}$
  and $S^1_{p_j}$ are glued to produce a node, raising the genus of
  the curve by~$1$ (non-separating) or connecting two components
  (separating).
\end{enumerate}
The factorization category $\cC$ is the \emph{primitive datum}: it
is not an algebra, but a category.  Its objects are boundary
conditions; its morphisms are open-string states.
\end{definition}

\begin{theorem}[Boundary algebra from vacuum;
  \ClaimStatusProvedHere]
\label{thm:boundary-algebra-from-vacuum-v2}
Let $\cC$ be an open/closed factorization dg-category on
$(X, D, \tau)$, and let $b \in \cC(I_p)$ be a compact generator
(\emph{vacuum object}).  Then:
\begin{enumerate}[label=\textup{(\roman*)}]
\item The endomorphism algebra
  $A_b := \RHom_{\cC(I_p)}(b, b)$
  is an $A_\infty$-algebra, with operations
  $\mu_n^{A_b}$ from ordered boundary collisions on
  $\Conf_n^{\mathrm{ord}}(I_p)$.

\item The mixed operations endow $A_b$ with the structure of a
  curved $A_\infty$-chiral algebra: the operations
  \[
    m_k \colon A_b^{\otimes k} \to
    A_b((\lambda_1))\cdots((\lambda_{k-1}))
  \]
  are obtained by restricting the bulk-boundary mixed operations
  $\mu_{k;0}$ to a formal disk at~$p$ and recording the collision
  data in spectral variables $\lambda_j = z_j - z_k$.

\item The $A_\infty$-chiral structure depends on the choice of
  compact generator~$b$.  Different choices
  $b, b' \in \cC(I_p)$ give Morita-equivalent $A_\infty$-chiral
  algebras:
  $\Perf(A_b) \simeq \Perf(A_{b'}) \simeq \cC(I_p)$.
  The Morita class of $A_b$ is an invariant of $\cC$; the specific
  algebra $A_b$ is a \emph{chart}.
\end{enumerate}
\end{theorem}

\begin{proof}
(i) The ordered boundary configuration space
$\Conf_n^{\mathrm{ord}}(I_p)$ is the interior of the simplex
$\Delta^{n-1}$.  Its Stasheff compactification is the associahedron
$K_n$, whose codimension-one faces parametrize consecutive-block
compositions.  Boundary corrections on $\partial K_n$ give the
$A_\infty$-relations $\sum_{i+j=n+1} \mu_i \circ \mu_j = 0$.

(ii) The formal disk $\widehat{D}_p$ at~$p$ identifies the bulk
configuration space with $\FM_k(\widehat{D}_p)$.  After choosing a
local coordinate $z$ centered at~$p$, the relative positions
$\lambda_j = z_j - z_k$ are the spectral variables.  The FM
compactification on the formal disk gives operadic compositions in
the chiral endomorphism operad $\cE\!nd^{\mathrm{ch}}_{A_b}$
(Definition~\ref{def:chiral-endomorphism-operad-v2}): codimension-one
faces correspond to operadic insertions with block-substitution
of spectral variables.

(iii) Morita invariance: if $b' \in \cC(I_p)$ is another compact
generator, the bimodule $\RHom(b, b')$ implements a Morita
equivalence between $A_b$ and $A_{b'}$.  The derived categories
$\Perf(A_b) \simeq \Perf(A_{b'})$ are canonically equivalent,
and both recover $\cC(I_p)$.
\end{proof}

\subsection{The universal bulk}

\begin{definition}[Chiral endomorphism operad]
\label{def:chiral-endomorphism-operad-v2}
Let $A$ be a graded vector space with translation
$\partial \colon A \to A$.  The \emph{chiral endomorphism operad}
$\cE\!nd^{\mathrm{ch}}_A$ has components
$\cE\!nd^{\mathrm{ch}}_A(n)
:= \Hom(A^{\otimes n},\,
A((\lambda_1))\cdots((\lambda_{n-1})))$
with composition by consecutive block-substitution of spectral
variables.
\end{definition}

\begin{definition}[Chiral Hochschild cochains]
\label{def:chiral-hochschild-v2}
The \emph{chiral Hochschild cochain complex} of a curved
$A_\infty$-chiral algebra $(A, \{m_k\})$ is
\[
  \cC^\bullet_{\mathrm{ch}}(A, A)
  \;:=\;
  \prod_{n \ge 0}
  \cE\!nd^{\mathrm{ch}}_A(n+1)[-n],
\]
with differential $\delta f := [m, f]$ using the Gerstenhaber
bracket from operadic braces.  The brace operations
$f\{g_1, \ldots, g_r\}$ are defined by iterated operadic
insertion with Koszul signs.
\end{definition}

\begin{theorem}[The universal bulk;
  \ClaimStatusProvedHere]
\label{thm:universal-bulk}
\index{derived center!as universal bulk|textbf}
\index{Swiss-cheese theorem!chiral|textbf}
Let $(A, \{m_k\})$ be a curved $A_\infty$-chiral algebra.
\begin{enumerate}[label=\textup{(\roman*)}]
\item \emph{Brace dg algebra.}  The complex
  $(\cC^\bullet_{\mathrm{ch}}(A, A), \delta, \{-\}\{-,\ldots,-\})$
  is a brace dg algebra: $\delta^2 = 0$ from $[m,m] = 0$, the
  brace operations satisfy the higher pre-Lie identity, and the
  Gerstenhaber bracket $[f,g] = f\{g\} - (-1)^{\|f\|\|g\|} g\{f\}$
  descends to a Gerstenhaber algebra on cohomology.

\item \emph{Swiss-cheese universality.}
  The pair $\cU(A) := (\cC^\bullet_{\mathrm{ch}}(A, A),\; A)$
  is the \emph{terminal} local chiral open/closed pair: for every
  local open/closed pair $(\cB, A, \iota)$, there exists a unique
  morphism $\Phi \colon \cB \to \cC^\bullet_{\mathrm{ch}}(A,A)$
  intertwining the Swiss-cheese operations.

\item \emph{Derived center.}  The \emph{chiral derived center}
  \[
    \cZ^{\mathrm{der}}_{\mathrm{ch}}(A)
    \;:=\;
    H^*(\cC^\bullet_{\mathrm{ch}}(A, A),\, \delta)
  \]
  is the universal bulk algebra.  On the Koszul locus, it is a
  three-term Gerstenhaber algebra
  $\cZ^0 \oplus \cZ^1 \oplus \cZ^2$
  by Theorem~H of Volume~I.

\item \emph{Morita invariance.}
  The derived center depends only on the Morita class of the
  open-sector category:
  $\cZ^{\mathrm{der}}_{\mathrm{ch}}(\cC)
  := R\!\Hom_{\operatorname{Fun}(\cC,\cC)}(\Id, \Id)$,
  which in any chart $\cC \simeq \Perf(A_b)$ is computed by
  $\cC^\bullet_{\mathrm{ch}}(A_b, A_b)$.
\end{enumerate}
\end{theorem}

The proof of each part is given in
\S\ref{sec:universal-bulk-proofs}.  The
derived center is the \emph{bulk}, not the bar complex.
The bar complex $\barB(A)$ classifies \emph{twisting morphisms}
--- universal couplings between $A$ and its Koszul dual~$A^!$
(Volume~I, Theorem~A) --- while the derived center classifies
\emph{bulk operators acting on the boundary}.  These are different
objects solving different problems.

\subsection{The bar complex as coalgebraic projection}

The bar construction of Volume~I enters the open/closed picture as
a \emph{derived invariant} of the factorization category~$\cC$, not
as the category itself.

\begin{proposition}[Bar complex from the open sector;
  \ClaimStatusProvedHere]
\label{prop:bar-from-open}
Let $\cC$ be an open/closed factorization dg-category with vacuum
object~$b$ and boundary algebra $A_b = \End(b)$.  The bar complex
$\barB^{\mathrm{ch}}(A_b) = T^c(s^{-1}\bar{A}_b)$ records the
twisting data of~$\cC$:
\begin{enumerate}[label=\textup{(\roman*)}]
\item The bar differential $d_{\barB}$ encodes how closed-sector
  operations twist the open-sector composition (OPE residues
  on $\FM_k(\C)$).

\item The bar coproduct~$\Delta$ encodes the ordered splitting
  of open-sector tensor factors (deconcatenation on
  $\Conf_k(\R)$).

\item Together, $(\barB(A_b), d_{\barB}, \Delta)$ is a dg coalgebra
  over the closed color of $\SCchtop$.  This is the Swiss-cheese
  fingerprint identified in Volume~I's Theorem~A.

\item The bar-cobar adjunction $\barB \dashv \Omega$ classifies
  twisting morphisms: $\Hom(\Omega(C), A_b) \cong
  \Tw(C, A_b) \cong \Hom(C, \barB(A_b))$.

\item The Koszul dual~$A^!$ is the algebra determined by the Verdier
  intertwining $D_{\Ran}(\barB(A)) \simeq \barB(A^!)$ (Theorem~A):
  it is the algebra whose bar coalgebra is the Verdier dual of~$\barB(A)$.
  It lives on the boundary~$\R$, governing the line-operator category
  $\cC_{\mathrm{line}} \simeq A^!\text{-}\mathsf{mod}$.
\end{enumerate}
Four functors, four outputs:
\begin{center}
\renewcommand{\arraystretch}{1.2}
\begin{tabular}{lll}
\textbf{Functor} & \textbf{Output} & \textbf{Role} \\
\hline
$\barB$ & coalgebra $\barB(A)$ & twisting morphism classifier \\
$\Omega$ & algebra $\Omega(\barB(A)) \simeq A$ & boundary recovery (inversion) \\
$D_{\Ran}$ & coalgebra $D_{\Ran}(\barB(A)) \simeq \barB(A^!)$ & Verdier intertwining \\
$\cC^\bullet_{\mathrm{ch}}(-,-)$
  & Gerstenhaber algebra $\cZ^{\mathrm{der}}$
  & universal bulk \\
\end{tabular}
\end{center}
None of these four objects should be confused with any other.
\end{proposition}

% ===================================================================
\section{Modularity}
\label{sec:modularity-trace-clutching}

Modularity is not an axiom on the closed algebra.
It is a consequence of three structures on the open sector.

\subsection{The annulus trace}

\begin{theorem}[Annulus trace; \ClaimStatusProvedHere]
\label{thm:annulus-trace-v2}
Let $\cC$ be an open/closed factorization dg-category on
$(X, D, \tau)$ with cyclic trace $\Tr_b$ at a puncture~$p$.
The factorization homology of the open sector over the boundary
circle $S^1_p$ recovers Hochschild homology:
\[
  \int_{S^1_p} \cC
  \;\simeq\;
  \HH_*(\cC).
\]
In a chart $\cC \simeq \Perf(A_b)$:
$\HH_*(A_b) = H_*(A_b \otimes_{A_b^e} A_b)$.
\end{theorem}

This is the \emph{first modular shadow}: the genus-$0$ trace
functional on the open sector.  It precedes the genus-$1$
curvature $\kappa(\cA) \cdot \lambda_1$ in the modular hierarchy:
the annulus is a trace, the torus is a self-trace (a trace composed
with itself along one cycle).  The passage from annulus to torus is
the passage from Hochschild homology to the genus-$1$ amplitude ---
and this passage is implemented by clutching.

\subsection{Clutching and the genus tower}

\begin{construction}[The clutching map]
\label{constr:clutching}
The non-separating clutching map
$\xi \colon \overline{\cM}_{g-1,n+2} \to \overline{\cM}_{g,n}$
(identifying two marked points to create a node, raising the genus
by~$1$) induces a sewing functor on the open sector:
\[
  \xi_* \colon
  \cC(\Sigma_{g-1,n+2})
  \;\longrightarrow\;
  \cC(\Sigma_{g,n}).
\]
The genus-$g$ amplitude is obtained by composing $g$~clutching maps,
starting from the genus-$0$ seed.  The separating clutching
$\xi_{\mathrm{sep}} \colon
\overline{\cM}_{g_1,n_1+1} \times \overline{\cM}_{g_2,n_2+1}
\to \overline{\cM}_{g,n}$ glues two surfaces along a shared
boundary circle.

Both clutching maps are implemented by the cyclic trace: sewing
along a boundary circle is the composition
\[
  \cC(\Sigma_1) \otimes \cC(\Sigma_2)
  \xrightarrow{\text{restrict to shared circle}}
  \cC(S^1)
  \xrightarrow{\Tr}
  k,
\]
where the trace on $\cC(S^1) \simeq \HH_*(\cC)$ extracts the
genus-raising scalar.
\end{construction}

\begin{theorem}[The modular Maurer--Cartan equation;
  \ClaimStatusProvedHere]
\label{thm:modular-mc}
\index{Maurer--Cartan equation!modular|textbf}
The open/closed MC element
$\Theta^{\mathrm{oc}} :=
\Theta_\cA + \sum_j \mu^{\cM_j}$
satisfies:
\begin{equation}\label{eq:modular-mc-v2}
  D^{\mathrm{oc}}\,\Theta^{\mathrm{oc}}
  + \tfrac{1}{2}[\Theta^{\mathrm{oc}},\Theta^{\mathrm{oc}}]
  + \hbar\,\Delta_{\mathrm{clutch}}(\Theta^{\mathrm{oc}})
  = 0.
\end{equation}
The three terms are:
\begin{enumerate}[label=\textup{(\alph*)}]
\item $D^{\mathrm{oc}}\,\Theta^{\mathrm{oc}}$: the differential
  (Stokes' theorem on the interior of the bordered FM
  compactification).

\item $[\Theta^{\mathrm{oc}},\Theta^{\mathrm{oc}}]$: the bracket
  (codimension-$1$ collision strata --- bubble trees).

\item $\hbar\,\Delta_{\mathrm{clutch}}(\Theta^{\mathrm{oc}})$:
  the clutching operator (codimension-$1$ nodal strata ---
  genus-raising degenerations), weighted by $\hbar$ to track genus.
\end{enumerate}
The equation is the boundary-of-boundary-is-zero identity on the
compactified bordered moduli spaces.  It decomposes by sector:
pure closed, pure open, mixed (bulk-to-boundary coupling), and
clutching (genus tower).  The pure-closed projection recovers
$\Theta_\cA$ (Volume~I, Theorem~\ref*{V1-thm:mc2-bar-intrinsic}).
\end{theorem}

\subsection{The modular trace principle}

\begin{principle}[Modularity $=$ trace $+$ clutching]
\label{princ:modularity-v2}
\index{modular trace principle|textbf}
The modular structure of the theory is not an additional axiom.
It is determined by three ingredients on the open sector:
\begin{center}
\renewcommand{\arraystretch}{1.3}
\begin{tabular}{lll}
\textbf{Ingredient} & \textbf{Open-sector origin}
  & \textbf{Genus contribution} \\
\hline
Cyclic trace & $\Tr_b \colon \End(b) \to k$
  & seed: the genus-$0$ amplitude \\
Annulus & $\int_{S^1} \cC \simeq \HH_*(\cC)$
  & genus $\frac{1}{2}$: the open partition function \\
Clutching & $\xi_* \colon \cC_{g-1} \to \cC_g$
  & each application raises genus by~$1$ \\
\end{tabular}
\end{center}
The genus-$g$ amplitude is the $g$-fold iterate of clutching
applied to the genus-$0$ seed.  The curvature
$\kappa(\cA) \cdot \omega_g$ (Volume~I, Theorem~D) is the leading
scalar of this iteration.  The shadow Postnikov tower
$(\kappa, \mathfrak{C}, \mathfrak{Q}, \ldots)$ is the arity
filtration.  The genus expansion is the genus filtration.  The full
MC element $\Theta^{\mathrm{oc}}$ encodes both simultaneously.
\end{principle}

% ===================================================================
\section{The gravitational interpretation}
\label{sec:gravitational-interpretation}

The preceding algebraic constructions have a single physical
content: they compute the gravitational partition function of
three-dimensional holomorphic--topological quantum field theory.

\subsection{The genus expansion as gravity}

The Maurer--Cartan element $\Theta^{\mathrm{oc}}$ is a formal
power series in the genus parameter~$\hbar$:
\[
  \Theta^{\mathrm{oc}}
  = \sum_{g \ge 0} \hbar^g \,\Theta^{\mathrm{oc}}_g.
\]
Each genus-$g$ component $\Theta^{\mathrm{oc}}_g$ lives in the
open/closed convolution algebra
$\mathfrak{g}^{\mathrm{oc}}_{\cA, \vec{\cM}}$
restricted to genus~$g$.  The genus-$0$ term is the tree-level
data (classical); the genus-$1$ term carries the curvature
$\kappa(\cA) \cdot \lambda_1$ (one-loop correction); each higher
genus carries the gravitational quantum corrections.

In the holomorphic--topological interpretation:
\begin{enumerate}[label=\textup{(\alph*)}]
\item \emph{Genus $0$}: the tree-level scattering.  The binary
  collision residue $r(z) = \Res^{\mathrm{coll}}_{0,2}(\Theta_\cA)$
  is the classical $r$-matrix, satisfying the classical Yang--Baxter
  equation by Stokes' theorem on $\FM_3(\C)$.  This is the leading
  gravitational vertex.

\item \emph{Genus $1$}: the one-loop correction.  The curvature
  $d_{\barB}^2 = \kappa(\cA) \cdot \omega_1$ is the Hodge anomaly:
  the genus-$1$ moduli space has a nontrivial Hodge bundle, and the
  fibrewise bar differential acquires curvature proportional to the first Chern
  class.  The modular characteristic $\kappa(\cA)$ is the
  gravitational anomaly coefficient --- it controls the strength
  of one-loop gravity.

\item \emph{Genus $g \ge 2$}: the higher-loop corrections.  The
  shadow tower $(\kappa, \mathfrak{C}, \mathfrak{Q}, \ldots)$
  extracts the gravitational observables at each arity.  The
  complementarity theorem (Volume~I, Theorem~C) gives:
  \[
    Q_g(\cA) \oplus Q_g(\cA^!)
    \;\cong\;
    H^*(\overline{\cM}_g,\, Z(\cA)),
  \]
  decomposing the moduli-space cohomology into complementary
  Lagrangian contributions from the algebra and its dual.
  This is the algebraic expression of the holographic principle:
  boundary and bulk contributions to the gravitational amplitude
  are complementary halves of a single symplectic vector space.
\end{enumerate}

\subsection{Boundary-holographic complexity}

The bulk 3d HT theory is topological: it has no local propagating
degrees of freedom.  The perturbative gravitational partition
function is determined entirely by the boundary chiral algebra.
The shadow depth $r_{\max}(\cA)$ classifies the nonlinear
complexity of this boundary-determined partition function
(Theorem~\ref{thm:thqg-boundary-holographic-complexity}).

\begin{center}
\renewcommand{\arraystretch}{1.3}
\begin{tabular}{clll}
\textbf{Class} & \textbf{Depth} & \textbf{Examples}
  & \textbf{Boundary-holographic content} \\
\hline
$\mathsf{G}$ & $r_{\max} = 2$ & Heisenberg, lattice
  & Gaussian: $F_g = \kappa \cdot \lambda_g^{\mathrm{FP}}$,
    one-loop exact \\
$\mathsf{L}$ & $r_{\max} = 3$ & affine $\widehat{\fg}_k$
  & Lie/tree: cubic boundary vertices only \\
$\mathsf{C}$ & $r_{\max} = 4$ & $\beta\gamma$
  & Contact: quartic terminates by rank-one rigidity \\
$\mathsf{M}$ & $r_{\max} = \infty$ & Virasoro, $\cW_N$
  & Mixed: infinite tower, non-polynomial potential \\
\end{tabular}
\end{center}

For class-$\mathsf{G}$ algebras, the boundary OPE is quadratic:
the partition function is determined at one loop.  Class~$\mathsf{L}$
has tree-level cubic boundary interactions but no higher.
Class~$\mathsf{C}$ terminates at quartic by a stratum-separation
mechanism.  Class~$\mathsf{M}$ has an infinite tower of
boundary $A_\infty$ operations, and the shadow generating function
$H(t)^2 = t^4 Q_L(t)$ is algebraic of degree~$2$ with
growth rate $\rho(\cA) = \sqrt{9\alpha^2 + 2\Delta}/(2|\kappa|)$
(Theorem~\ref*{V1-thm:shadow-radius}).  The Virasoro algebra at
$c = 13$ (the self-dual point, $\kappa = \kappa^! = 13/2$) has
$\rho \approx 0.467$: the gravitational tower converges.

\subsection{The completed platonic datum}

The entire structure is packaged in a single datum.

\begin{definition}[Completed platonic datum]
\label{def:completed-platonic-v2}
\index{platonic datum!completed|textbf}
For a gravitational input $(\cA, X, \langle\cdot,\cdot\rangle, k)$
on a tangential log curve $(X, D, \tau)$ with boundary category
$\cC$, the \emph{completed platonic datum} is:
\[
  \Pi^{\mathrm{oc}}_X(\cA)
  \;:=\;
  \bigl(\,
  \underset{\text{boundary}}{\cA},\;\;
  \underset{\text{lines}}{A^!},\;\;
  \underset{\text{bulk}}{\cZ^{\mathrm{der}}_{\mathrm{ch}}(\cC)},\;\;
  \underset{\text{MC element}}{\Theta^{\mathrm{oc}}},\;\;
  \underset{\text{$r$-matrix}}{r(z)},\;\;
  \underset{\text{connection}}{\nabla^{\mathrm{hol}}},\;\;
  \underset{\text{trace}}{\Tr_\cA},\;\;
  \underset{\text{tower}}{\mathfrak{R}^{\mathrm{oc}}_\bullet}
  \,\bigr).
\]
Boundary, bulk, and lines are independently defined but constrained
by the open/closed MC equation.  The closed projection of
$\Theta^{\mathrm{oc}}$ is the modular MC element $\Theta_\cA$ of
Volume~I.  The mixed projection is the bulk-to-boundary coupling.
The clutching projection is the genus tower.  Every gravitational
observable is a component of this single object.
\end{definition}
